\section{Discussion}

The replication supports two distinct conclusions.  First, a decreasing Biot
coefficient does not require the coefficient itself to be an independently
fitted pressure function.  It follows from a pressure-dependent constrained
tangent: cracks close, the drained skeleton stiffens toward the unjacketed
solid response, and the ratio \(K/K_s'\) increases.  The same physical trend is
represented by the finite-deformation derivative in
\cref{eq:nonlinear-biot}.

Second, agreement with Figure~3c alone does not uniquely validate the selected
crack law.  The reported coefficient is calculated from the measured moduli, so
the present comparison primarily tests whether a compact, positive constitutive
model can reproduce those moduli and whether the MOOSE implicit derivative
returns the resulting coefficient consistently.  Independent validation would
require volumetric strain, pore-volume change, or intrinsic density measurements
under controlled pore pressure.  The finite-deformation definition suggests a
particularly direct experiment: measure the intrinsic solid specific-volume
change as the skeleton volume changes while pore pressure is held fixed.

The SP14 workbook discrepancy illustrates the importance of archiving source
cells and formulas.  Using the annotated \(65.5\)~GPa rather than the workbook's
65.0~GPa shifts every SP14 coefficient.  The present repository therefore
preserves the original workbook, checksums, source rows, recomputed values, and
the exact modulus used by the released formulas.

Reaction is absent from the present calculation.  This separation is deliberate:
HT14, SP14, and SP30 are treated as three fixed materials rather than states on
a simulated reaction trajectory.  A subsequent reacting study can allow the
free-energy, crack-compliance, and unjacketed-modulus parameters to evolve with
registered phase/component reaction variables.  In that case the tangent must
state whether reaction/internal variables are held fixed or included through
the local implicit system.

